\begin{enumerate}
	\item Le triangle OAD est rectangle en A.

	D'après le théorème de Pythagore, on a :

	$\begin{aligned}
	\mathrm{OD}^2&=\mathrm{AD}^2+\mathrm{AO}^2\\
	8,2^2&=1,8^2+\mathrm{AO}^2\\
	\mathrm{AO}^2&=8,2^2-1,8^2\\
	\mathrm{AO}^2&=64\\
	\mathrm{AO}&=\sqrt{64}& \text{car une longueur est positive.}\\
	\mathrm{AO}&=8.
	\end{aligned}$

	Donc la longueur du segment [AO] est bien de 8 cm.

	\item (BC) et (AD) sont toutes les deux perpendiculaires à la même droite (BO).

	Donc les droites (BC) et (AD) sont donc parallèles.

	\item \begin{itemize}
		\item Les points O, A, B dans cet ordre;
		\item les points O, D, C sont alignés, dans le même ordre;
		\item les droites (BC) et (AD) sont parallèles.
	\end{itemize}

	D'après le théorème de Thalès, on a :\quad $\dfrac{\mathrm{OA}}{\mathrm{OB}}=\dfrac{\mathrm{AD}}{\mathrm{BC}}=\dfrac{\mathrm{AD}}{\mathrm{BC}}$.

	Notamment : \quad $\dfrac{\mathrm{OA}}{\mathrm{OB}}=\dfrac{\mathrm{AD}}{\mathrm{BC}}$.

	En remplaçant par les valeurs connues :\quad $\dfrac{8}{\mathrm{OB}}=\dfrac{1,8}{4,5}$

	Finalement,\quad : $\mathrm{OB}=\dfrac{4,5\times 8}{1,8} =20$.

	Donc la longueur du segment [OB] est 20 cm.

	\item \begin{enumerate}
		\item Le volume du grand cône de hauteur [OB] se calcule ainsi :

		$V_G=\dfrac{\pi \times 4,5^2\times 20}{3}=135\pi\approx \np[cm^3]{424}$.

		\item Calculons le volume du petit cône de hauteur [AO] :

		$V_p=\dfrac{\pi \times 1,8^2\times 8}{3}=\dfrac{216\pi}{25}\approx \np[cm^3]{27}$

		Le volume du gobelet est donc :

		$V_G-V_p\approx 424-21=\np[cm^3]{397}$.
	\end{enumerate}
\end{enumerate}
